\section{训练与推理流程}

\subsection{训练算法}

Conditional Flow Matching 的完整训练算法如下：

\begin{lstlisting}[caption={Conditional Flow Matching 训练算法（OT 路径）}]
# v_theta: neural network parameterizing the vector field
# p_data: data distribution (training dataset)

for each training step:
    # 1. Sample data and noise
    x1 = sample_from(p_data)           # data sample
    x0 = sample_from(Normal(0, I))     # noise sample
    t  = sample_from(Uniform(0, 1))    # time step

    # 2. Compute interpolation (conditional flow map)
    xt = (1 - t) * x0 + t * x1

    # 3. Compute target (conditional vector field)
    target = x1 - x0

    # 4. Compute loss and update
    prediction = v_theta(t, xt)
    loss = MSE(prediction, target)
    loss.backward()
    optimizer.step()
\end{lstlisting}

\begin{warningbox}{训练中的注意事项}
\begin{enumerate}
    \item 时间 $t$ 应均匀采样自 $[0, 1]$，某些实现会对端点做小幅裁剪（如 $[\epsilon, 1-\epsilon]$）以避免数值问题
    \item 对于 OT 路径，目标 $x_1 - x_0$ 不依赖于 $t$，但 $x_t$ 依赖于 $t$，因此网络仍需以 $t$ 为输入
    \item 网络架构可以复用扩散模型的 U-Net 或 DiT 架构，只需确保时间 $t$ 的编码方式正确
\end{enumerate}
\end{warningbox}

\subsection{推理（采样）算法}

训练完成后，生成新样本的过程如下：

\begin{lstlisting}[caption={Flow Matching 采样算法（Euler 方法）}]
# v_theta: trained vector field network
# N: number of discretization steps
# dt: step size = 1/N

# 1. Sample from reference distribution
x = sample_from(Normal(0, I))

# 2. Solve ODE from t=0 to t=1
for i in range(N):
    t = i / N
    x = x + v_theta(t, x) * dt    # Euler step

# x is now approximately a sample from p_data
return x
\end{lstlisting}

可以使用更高阶的 ODE 求解器（如 Heun、Dormand-Prince/RK45）来提高采样质量或减少步数。

\subsection{本章小结}

Flow Matching 的训练和推理流程都极其简洁：训练就是在随机插值点上回归速度目标，推理就是从噪声出发求解 ODE。整个流程不涉及任何 SDE 模拟或复杂的噪声调度设计。


